\documentclass{ximera}
% xmPreamble.tex en xmPrintstyle.sty worden automatisch geladen door ximera.cls (v2.7.8);
% NIET nog eens \input{./preamble.tex} of \addPrintStyle{.} doen (geeft 'already defined'-fouten).

\makeatletter % eigen img-map toevoegen aan het zoekpad (werkt standalone én vanuit de xourse)
\g@addto@macro\Ginput@path{{img/}{UW4202/img/}}
\makeatother

\begin{document}
\author{Alexander Holvoet}
\xmtitle{Zand op een schijf met een gat in}{}

Het grondvlak waarop we zand willen gieten, is een cirkelschijf met een gaatje in. Je ziet
hieronder twee cirkelschijven: \'e\'en met een gat in het midden, \'e\'en met een
gedecentreerd gat. Het gat is klein en in de redeneringen behandelen we het als een punt.
In werkelijkheid mag het gat ook niet \emph{te} klein zijn, want de zandkorrels moeten
erdoor kunnen.

\begin{image}[12cm]
    \includegraphics[width=5.5cm,keepaspectratio]{zand07.jpg}
    \quad
    \includegraphics[width=5.5cm,keepaspectratio]{zand08.jpg}
\end{image}
Links een centraal gat, rechts een gedecentreerd gat.

\begin{exercise}
    Teken op deze bovenaanzichten telkens de kamlijn, giet dan zand op het kartonnen
    exemplaar om dit te verifi\"eren, en verklaar ten slotte de vorm van de kamlijn.
    Begin met het \textbf{centrale} gat.
    \begin{hint}
        De kamlijn bestaat uit de punten die even ver liggen van het gat als van de rand
        van de schijf. Welke kromme krijg je bij een centraal gat?
    \end{hint}
    \begin{oplossing}
        Bij het centrale gat is het bovenaanzicht van de kamlijn een concentrische cirkel
        met als straal de helft van de straal van de gegeven schijf. Deze cirkel bestaat
        immers uit de punten op gelijke afstand tot het gat en de rand. In 3D heb je in het
        midden van de kegel een kegelvormige uitsparing met dezelfde rusthoek: de zandhoop
        is een kegel waaruit een tweede, omgekeerde kegel met dezelfde rusthoek is
        weggehaald. Die tweede kegel bepaalt de uitsparing.
        \begin{image}[12cm]
            \includegraphics[width=4cm,keepaspectratio]{zand09.jpg}
            \quad
            \includegraphics[width=6cm,keepaspectratio]{zand10.jpg}
        \end{image}
        Links het boven- en vooraanzicht, rechts de kegel met de omgekeerde kegel die de
        uitsparing bepaalt.
    \end{oplossing}
\end{exercise}

\begin{exercise}
    Bekijk nu het \textbf{gedecentreerde} gat. Welke kromme is de kamlijn hier in
    bovenaanzicht? Bewijs je vermoeden.
    \begin{hint}
        Een punt \(P\) van de kamlijn ligt even ver van het gat \(G\) als van de rand van
        de schijf. Meet de afstand tot de rand loodrecht, dus langs een straal. Wat kun je
        zeggen over \(|PM| + |PG|\), met \(M\) het middelpunt van de schijf?
    \end{hint}
    \begin{oplossing}
        De kamlijn is in bovenaanzicht een \emph{ellips}. Als het gat dicht bij het
        middelpunt ligt, lijkt het op het eerste gezicht een cirkel, maar laten we
        redeneren. Een willekeurig punt \(P\) van de kamlijn ligt even ver van het gat \(G\)
        als van de rand van de schijf. Die afstand \(|PQ|\) tot de rand wordt loodrecht op
        die rand gemeten, dus langs een straal. Bijgevolg is
        \[
            |PM| + |PG| = |PM| + |PQ| = |MQ| = r ,
        \]
        waarbij \(r\) de straal van de cirkel is. De som \(|PM| + |PG|\) is dus constant, en
        de punten \(P\) van de kamlijn vormen een ellips met als brandpunten \(G\) en \(M\).
        \begin{image}[4cm]
            \includegraphics[width=4cm,keepaspectratio]{zand11.jpg}
        \end{image}
    \end{oplossing}
\end{exercise}

\begin{exercise}
    Je kunt, bij het gedecentreerde gat, ook redeneren met hoogtelijnen. Je ziet hier de
    cirkelschijf met het gat, en concentrische cirkels rond het middelpunt en rond het gat.
    Construeer met deze figuur de kamlijn en leg uit.
    \begin{image}[4cm]
        \includegraphics[width=4cm,keepaspectratio]{zand12.jpg}
    \end{image}
    \begin{oplossing}
        Je kunt in het vlak redeneren: de punten van de kamlijn liggen op gelijke afstand
        van het gat als van de rand. Bijvoorbeeld 5 eenheden van het gat en 5 eenheden van
        de rand, waarbij een eenheid de afstand tussen twee concentrische cirkels is. Dit
        geeft de punten \(P\) en \(P'\). Analoog voor andere aantallen eenheden.

        Je kunt ook in 3D redeneren: de cirkels zijn hoogtelijnen, en je snijdt de cirkels
        die even hoog zijn. De cirkels rond het gat beginnen in het gat met hoogte nul
        (zoals de kegel op zijn top); de cirkels rond het middelpunt beginnen aan de rand
        met hoogte nul (zoals de gewone kegel).
        \begin{image}[4cm]
            \includegraphics[width=4cm,keepaspectratio]{zand13.jpg}
        \end{image}
    \end{oplossing}
\end{exercise}

\begin{exercise}
    Het gat was tot nu toe heel klein en we behandelden het als een punt. Wat verandert er
    als het (gedecentreerde) gat een grotere cirkel(schijf) is?
    \begin{oplossing}
        Het bovenaanzicht van de kamlijn is opnieuw een ellips. De brandpunten zijn het
        middelpunt \(M\) van de grote cirkel en het middelpunt \(G\) van het gat. De som van
        de afstanden tot deze twee punten is weer constant, maar deze constante is nu niet
        de straal van de cirkel, wel de som van de straal van de cirkel en die van het gat.

        Een punt \(P\) van de kamlijn ligt even ver van de rand van de grote cirkel als van
        de rand van het ronde gat: \(|PQ| = |PH|\). De som van de afstanden van \(P\) tot
        \(M\) en \(G\) is bijgevolg
        \[
            |PM| + |PG| = |PM| + |PH| + |HG| = |PM| + |PQ| + |HG| = |MQ| + |HG| = r + s ,
        \]
        waarbij \(r\) de straal van de grote cirkel is en \(s\) de straal van het gat. Deze
        som is constant, dus \(P\) beschrijft een ellips.
        \begin{image}[12cm]
            \includegraphics[width=4cm,keepaspectratio]{zand14.jpg}
            \quad
            \includegraphics[width=6cm,keepaspectratio]{zand15.jpg}
        \end{image}
    \end{oplossing}
\end{exercise}

\end{document}
